\documentclass[11pt,a4paper]{article}
\usepackage[T1]{fontenc}
\usepackage[utf8]{inputenc}
\usepackage[margin=20mm]{geometry}
\setlength{\parindent}{0pt}
\setlength{\parskip}{6pt}
\begin{document}
\begin{center}\large
Computer Science BSc, Discrete Mathematics I\\
Rules of earning plus points in the practice course\\
(@@SEMESTER@@)
\end{center}

Plus points can be earned during the study period by solving optional additional
problems set by the practice class teacher.
\textbf{Maximum plus points:} @@PLUS_POINTS_MAXIMUM@@

\section*{Use of plus points}
@@TOTAL_SCORE_RULES@@

\section*{Tasks that can be solved for plus points}
\textbf{Eligibility, scoring, deadlines and methods of handing in solutions:}\par
@@PLUS_POINTS_RULES@@

\subsection*{Homework set in classes}
Homework provides practice in the material covered. The homework questions
indicated by the practice class teacher can be submitted for plus points according
to the scoring, deadlines and submission rules above.

\subsection*{``Further questions'' in problem sets}
In the ``Further questions'' section you can find two types of questions:
questions that do not require coding and coding questions. For a full solution,
always prove why your answer is correct, except for questions that require writing
a program. Follow the separate submission rules for questions without coding and
coding questions given above.

\subsection*{Practice quizzes published on Canvas for plus points}
Practice quizzes on Canvas for plus points will be published during the semester
that can also be solved for plus points. The applicable quiz names, score
thresholds, awards and deadlines are specified in the rules above.

\subsection*{Other optional tasks set by the practice class teacher}
Your practice class teacher may set any other optional tasks which can be completed
during the semester for plus points. The scoring and deadlines of these will be
communicated to you in due course by your practice class teacher.
\end{document}
